\begin{center}
\SourceLeaf{046}\pdfbookmark[2]{30.}{article-030}30.
\end{center}
Quum omnes quos e tabulis logarithmicis et trigonometricis depromimus numeri praecisionem absolutam non admittant, sed ad certum tantummodo gradum sint approximati, ex omnibus calculis illarum adiumento perfectis proxime tantum vera resultare possunt. In plerisque quidem casibus tabulae vulgares ad septimam figuram decimalem usque exactae, i. e. ultra dimidiam unitatem in figura septima excessu seu defectu nunquam aberrantes a vero, praecisionem plus quam sufficientem suppeditant, ita ut errores inevitabiles nullius plane sint momenti: nihilominus utique fieri \`{p}otest, ut errores tabularum in casibus specialibus effectum suum exserant augmentatione tanta, ut methodum alias optimam plane abdicare aliamque ei substituere cogamur. Huiusmodi casus in iis quoque calculis, quos hactenus explicavimus, occurrere potest; quamobrem ab instituto nostro haud alienum erit, disquisitiones quasdam circa gradum praecisionis, quam tabulae vulgares in illis permittunt, hic instituere. Etsi vero ad hoc argumentum calculatori practico gravissimum exhauriendum hic non sit locus, investigationem eo perducemus, ut ad propositum nostrum sufficiat, et a quolibet, cuius interest, ulterius expoliri et ad quasvis alias operationes extendi possit.

\begin{center}
\hypertarget{native-046-article031}{}\pdfbookmark[2]{31.}{article-031}31.
\end{center}
Quilibet logarithmus, sinus, tangens etc. (aut generaliter quaelibet quantitas irrationalis e tabulis excerpta) errori obnoxius est, qui ad dimidiam unitatem in figura ultima ascendere potest: designabimus hunc erroris limitem per $\omega$, qui itaque in tabulis vulgaribus fit $=0{,}000\,0000\,5$. Quodsi logarithmus etc. e tabulis immediate desumi non potuit, sed per interpolationem erui debuit, error duplici caussa aliquantulum adhuc maior esse potest. \emph{Primo} enim pro parte proportionali, quoties (figuram ultimam tamquam unitatem spectando) non est integer, adoptari solet integer proxime maior vel minor: hac ratione errorem tantum non usque ad duplum augeri posse facile perspicitur. Ad hanc vero erroris augmentationem omnino hic non respicimus, quum nihil obstet, quominus unam alteramve figuram decimalem parti illi proportionali \SourceLeaf{047}affigamus, nulloque negotio pateat, logarithmum interpolatum, si pars proportionalis absolute exacta esset, errori maiori obnoxium non esse quam logarithmos in tabulis immediate expressos, quatenus quidem horum variationes tamquam uniformes considerare liceat. Erroris augmentatio \emph{altera} inde nascitur, quod suppositio ista omni rigore non est vera: sed hanc quoque negligimus, quoniam effectus differentiarum secundarum altiorumque in omnibus propemodum casibus nullius plane momenti est (praesertim si pro quantitatibus trigonometricis tabulae excellentissimae quas \textsc{Taylor} curavit adhibentur), facilique negotio ipsius ratio haberi possit, ubi forte paullo maior evaderet. Statuemus itaque pro omnibus casibus tabularum errorem maximum inevitabilem $=\omega$, siquidem argumentum (i. e. numerus cuius logarithmus, seu angulus cuius sinus etc. quaeritur) praecisione absoluta habetur. Si vero argumentum ipsum proxime tantum innotuit, errorique maximo, cui obnoxium esse potest, respondere supponitur logarithmi etc. variatio $\omega'$ (quam per rationem differentialium definire licet), error maximus logarithmi per tabulas computati usque ad $\omega+\omega'$ ascendere potest.

Vice versa, si adiumento tabularum argumentum logarithmo dato respondens computatur, error maximus ei eius variationi aequalis est, quae respondet variationi $\omega$ in logarithmo, si hic exacte datur, vel quae respondet variationi logarithmi $\omega+\omega'$, si logarithmus ipse usque ad $\omega'$ erroneus esse potest. Vix opus erit monere, $\omega$ et $\omega'$ eodem signo affici debere.

Si plures quantitates intra certos tantum limites exactae adduntur, aggregati error maximus aequalis erit aggregato singulorum errorum maximorum, iisdem signis affectorum; quare etiam in subtractione quantitatum proxime exactarum differentiae error maximus summae errorum singulorum maximorum aequalis erit. In multiplicatione vel divisione quantitatis non absolute exactae error maximus in eadem ratione augetur vel diminuitur ut quantitas ipsa.
